\section{数列习题}

\begin{example}
    数列 $\left\{a_n\right\}$ 的前 $n$ 项和为 $S_n$ ，则下列说法不正确的是 （\qquad）
    \begin{tasks}(2)
        \task 若 $a_n=-2 n+13$ ，则数列 $\left\{a_n\right\}$ 的前 5 项和 $S_5$ 最大
        \task 若等比数列 $\left\{a_n\right\}$ 是递减数列，则公比 $q$ 满足 $0<q<1$
        \task 在等差数列中，若 $S_{2025}>0$ ，则 $a_{1013}>0$
        \task 已知 $\left\{a_n\right\}$ 为等差数列，则数列 $\left\{\frac{S_n}{n}\right\}$ 也是等差数列
    \end{tasks}
\end{example}

\vspace{10em}

\begin{example}
    已知等差数列 $\left\{a_n\right\}$ 的前 $n$ 项和为 $S_n$ ，若 $a_1=29$ ，且 $S_{10}=S_{20}$ ，则 $S_n$ 的最大值为 \underline{\qquad}.
\end{example}

\v{8em}

\begin{example}
    数列 $\left\{a_n\right\}$ 的前 $n$ 项和为 $S_n$ ，已知 $S_n=-n^2+7 n$ ，则下列说法正确的是 （\qquad）
    \begin{tasks}(2)
        \task $\left\{a_n\right\}$ 是递增数列
        \task $a_{10}=-12$
        \task 当 $n>4$ 时，$a_n<0$
        \task 当 $n=5$ 时，$S_n$ 取得最大值
    \end{tasks}
\end{example}

\vspace{10em}

\begin{example}
    已知等差数列 $\left\{a_n\right\}$ 的首项为 20 ，公差为 -2 ，则数列 $\left\{a_n\right\}$ 的前 $n$ 项和 $S_n$ 的最大值为 （\qquad）
    \begin{tasks}(4)
        \task 100
        \task 108
        \task 112
        \task 110
    \end{tasks}
\end{example}

\vspace{10em}

\begin{example}
    设 $S_n$ 为等差数列 $\left\{a_n\right\}$ 的前 $n$ 项和，若 $a_1=7, S_3=15$ ．
    \begin{enumerate}
        \item 求数列 $\left\{a_n\right\}$ 的通项公式；

        \item 求数列 $\left\{a_n\right\}$ 的前 $n$ 项和 $S_n$ ，并求当 $n$ 为何值时，数列 $\left\{a_n\right\}$ 前 $n$ 和 $S_n$ 最大，求其最大值；

        \item 若数列 $\left\{b_n\right\}$ 的通项公式 $b_n=n$ ，求证：$\frac{1}{b_1 b_2}+\frac{1}{b_2 b_3}+\frac{1}{b_3 b_4}+\cdots+\frac{1}{b_n b_{n+1}}<1$
    \end{enumerate}
\end{example}

\vspace{10em}

\begin{example}
    等差数列 $\left\{a_n\right\}$ 的前 $n$ 项和为 $S_n, a_2 a_3=40, S_3=15$ ，数列 $\left\{b_n\right\}$ 满足 $\frac{1}{2} b_1+\frac{1}{2^2} b_2+\frac{1}{2^3} b_3+\cdots+\frac{1}{2^n} b_n=n\left(n \in \mathrm{~N}^*\right)$
    \begin{enumerate}
        \item 求数列 $\left\{a_n\right\}$ 和 $\left\{b_n\right\}$ 的通项公式；

        \item 若从数列 $\left\{a_n\right\}$ 中依次剔除与数列 $\left\{b_n\right\}$ 的公共项，剩下的项组成新的数列 $\left\{c_n\right\}$ ，求数列 $\left\{c_n\right\}$ 的前 50 项和 $T_{50}$ ．
    \end{enumerate}
\end{example}

\vspace{10em}

\begin{example}
    在数列 $\left\{a_n\right\}$ 中，$a_1=2, a_2=8$ ，且对任意的 $n \in \mathbf{N}^*$ ，都有 $a_{n+2}=4 a_{n+1}-4 a_n$ ，则 $\left\{a_n\right\}$ 的通项公式为 $\qquad$ $\therefore$ 若 $b_n=\left\{\begin{array}{l}\frac{n}{a_n}, n=2 k-1, k \in \mathbf{N}^* \\ \log _2 \frac{n}{a_n}, n=2 k, k \in \mathbf{N}^*\end{array}\right.$ ，则数列 $\left\{b_n\right\}$ 的前 $n$ 项和 $T_n=$ \underline{\qquad}．
\end{example}

\vspace{10em}

